\subsection{临界梳齿的 Cesaro--Haar 序参量、正性 Fredholm 完成障碍与 universal solenoid 的有限见证}\label{subsec:conclusion-rhsharp-comb-fredholm-solenoid-finite-witness}
前文已经分别固定了真实输入 \(40\) 状态核的 RH-sharp 谱界与隐藏的 \((-A_0)\) 因子、循环提升的临界线梳齿、通用 solenoid 上的有限素数见证、相位--Witt 转移算子的 Mellin 对角化、原子 Euler 乘积到 Fredholm--Witt 截面的构造性反向截面与谱刚性，以及 \(\Xi_\zeta\) 在圆盘完成化下的 Carath\'eodory/Toeplitz 正性等价链。把这些接口并读后，可再抽取出一组终端命题：临界梳齿并不只是“稠密化”，而可被视为一个具有 Haar 化极限的统计序参量；边界原子一旦进入正性 Fredholm 完成，便成为不能在完成层被抹除的 obstruction；对任意有限生成观测，这类 obstruction 的存在与否早在某个有限 prime-register 阶段已经被完全决定；而经典 RH 本身也可在 arithmetic solenoid 上被重写为“唯一幺正 Mellin 切片承载全部解析谱支撑”的命题。

\begin{theorem}[临界梳齿的 Cesaro--Haar 化定律]\label{thm:conclusion-critical-comb-cesaro-haar}
设 \(M\) 为有限核矩阵，\(\rho(M)=\lambda>1\)。设 \(\mu\neq\lambda\) 为一个边界模态，满足
\[
|\mu|=\sqrt{\lambda}.
\]
对每个整数 \(q\ge 2\)，令
\[
M^{\langle q\rangle}:=M\otimes P_q,
\]
其中 \(P_q\) 为 \(q\)-循环置换矩阵；再令
\[
\mathbb T_\lambda:=\RR/(2\pi/\log\lambda)\ZZ.
\]
把 \(\mu\) 在 \(M^{\langle q\rangle}\) 上诱导的临界线极点虚部写成 \(\mathbb T_\lambda\) 上的归一化计数测度
\[
\nu_q:=
\frac{1}{q}\sum_{j=0}^{q-1}
\delta_{-\frac{\arg\mu+2\pi j/q}{\log\lambda}},
\]
其中 Dirac 点取于 \(\mathbb T_\lambda\) 中对应的剩余类。则其 Cesaro 平均满足
\[
\frac1Q\sum_{q\le Q}\nu_q \rightharpoonup m_{\mathrm{Haar}}
\qquad (Q\to\infty),
\]
其中 \(m_{\mathrm{Haar}}\) 为 \(\mathbb T_\lambda\) 上的 Haar 概率测度。反之，若 \(M\) 满足严格核版 RH，即
\[
\max_{\nu\neq \lambda}|\nu|<\sqrt{\lambda},
\]
则不存在此类边界模态，因而临界梳齿测度族为空。
\end{theorem}
\leanverified{paper\_conclusion\_critical\_comb\_cesaro\_haar}
\begin{proof}
由命题 \ref{prop:finite-rh-phase-lift}，\(P_q\) 的谱为全部 \(q\) 次单位根 \(\omega_q^j=e^{2\pi i j/q}\)，故 \(\mu\) 在 \(M^{\langle q\rangle}\) 上裂解为
\[
\mu\omega_q^j,\qquad 0\le j\le q-1.
\]
在 Dirichlet 坐标 \(z=\lambda^{-s}\) 下，极点方程为
\[
\lambda^{-s}=(\mu\omega_q^j)^{-1}.
\]
取对数并模去周期 \(2\pi i/\log\lambda\)，得到相应极点满足
\[
\Re(s)=\frac{\log|\mu|}{\log\lambda}=\frac12,
\qquad
\Im(s)\equiv -\frac{\arg\mu+2\pi j/q}{\log\lambda}
\pmod{2\pi/\log\lambda}.
\]
这正给出上式中的测度 \(\nu_q\)。

对每个 \(n\in\ZZ\)，以字符 \(t\mapsto e^{in t\log\lambda}\) 定义 Fourier 系数，则
\[
\widehat{\nu_q}(n)
:=
\int_{\mathbb T_\lambda}e^{in t\log\lambda}\,\mathrm d\nu_q(t)
=
e^{-in\arg\mu}\cdot \frac1q\sum_{j=0}^{q-1}e^{-2\pi i nj/q}
=
e^{-in\arg\mu}\,\mathbf 1_{q\mid n}.
\]
当 \(n\neq 0\) 时，
\[
\left|
\frac1Q\sum_{q\le Q}\widehat{\nu_q}(n)
\right|
\le
\frac1Q\sum_{q\le Q}\mathbf 1_{q\mid n}
\le
\frac{d(|n|)}{Q}
\longrightarrow 0,
\]
其中 \(d(|n|)\) 为 \(n\) 的正因子个数；而当 \(n=0\) 时，上式恒等于 \(1\)。由紧阿贝尔群上 Haar 测度的 Fourier 刻画，即得所示弱收敛。

若严格核版 RH 成立，则由推论 \ref{cor:xi-critical-line-poles-empty-or-dense}，临界线 \(\Re(s)=\tfrac12\) 上不存在任何非平凡极点，故相应的临界梳齿测度族为空。证毕。
\end{proof}

\begin{theorem}[40 状态核的单通道饱和刚性]\label{thm:conclusion-realinput40-single-channel-comb-rigidity}
对真实输入 \(40\) 状态核 \(M\)，其非 Perron 边界模态恰有且仅有一个：
\[
\mu_\ast=-\sqrt{\lambda}=-\varphi,
\qquad
\lambda=\varphi^2.
\]
更强地，这一唯一达界模态完全由隐藏的 \((-A_0)\) 因子贡献。因而定理 \ref{thm:conclusion-critical-comb-cesaro-haar} 中的 Haar 化临界序参量在该核上只有一个非平凡通道。等价地，全部循环提升中的非平凡临界梳齿都由同一个符号扭曲黄金块生成，而其余纵向模态全部严格亚临界。
\end{theorem}
\begin{proof}
由定理 \ref{thm:pom-rh-sharp}，
\[
\mathrm{spec}(M)\setminus\{0\}
=
\{\lambda,\ -\sqrt{\lambda},\ 1^{(2)},\ (-1)^{(3)},\ \lambda^{-1},\ \lambda^{-1/2}\},
\qquad
\lambda=\varphi^2.
\]
因而除 Perron 根 \(\lambda\) 外，唯一满足 \(|\nu|=\sqrt{\lambda}\) 的非平凡谱点就是 \(-\sqrt{\lambda}\)。

另一方面，由命题 \ref{prop:real-input-40-A0-tensor-hidden}，
\[
\det(I-zM)
=(1-z)^2(1+z)\,
\det\!\bigl(I-z(A_0\otimes A_0)\bigr)\,
\det\!\bigl(I-z(-A_0)\bigr),
\]
其中
\[
A_0=
\begin{pmatrix}
1&1\\
1&0
\end{pmatrix}.
\]
由 \(\mathrm{spec}(A_0)=\{\varphi,-\varphi^{-1}\}\) 得
\[
\mathrm{spec}(A_0\otimes A_0)=\{\varphi^2,-1,-1,\varphi^{-2}\},
\qquad
\mathrm{spec}(-A_0)=\{-\varphi,\varphi^{-1}\}.
\]
于是达界点 \(-\sqrt{\lambda}=-\varphi\) 并非来自 \(40\) 状态复杂拼装的偶然共振，而是由 \(\det(I-z(-A_0))\) 强制出现。

故该核上的边界饱和通道唯一，且其来源就是单一的符号扭曲黄金因子。再结合定理 \ref{thm:conclusion-critical-comb-cesaro-haar}，便得到所述单通道 Haar 化结论。证毕。
\end{proof}

\begin{theorem}[循环相位细分的归一化自由能不可见性]\label{thm:conclusion-cyclic-lift-normalized-freeenergy-invisibility}
设 \(M\) 为有限维矩阵，\(\lambda=\rho(M)>1\)，并记 \(\{\mu_j\}_{j\in J}\) 为其全部非零特征值（按代数重数计）。对每个 \(q\ge 2\)，令
\[
M^{\langle q\rangle}:=M\otimes P_q,
\]
并定义归一化 Jensen--Mahler 自由能
\[
\mathcal G_{M^{\langle q\rangle}}(\sigma)
:=
\frac1q\int_0^{2\pi}
\log\left|\det\!\left(I-\lambda^{-\sigma}e^{i\theta}M^{\langle q\rangle}\right)\right|
\frac{d\theta}{2\pi}.
\]
则对一切 \(\sigma\in\RR\) 都有严格恒等式
\[
\mathcal G_{M^{\langle q\rangle}}(\sigma)
=
\int_0^{2\pi}
\log\left|\det\!\left(I-\lambda^{-\sigma}e^{i\theta}M\right)\right|
\frac{d\theta}{2\pi}
=
\sum_{j\in J}\max\!\bigl(0,\log|\mu_j|-\sigma\log\lambda\bigr).
\]
特别地，循环 lift 仅在临界线方向加密极点梳齿，而不改变归一化宏观自由能。若 \(M\) 在临界圆 \(|\mu|=\sqrt{\lambda}\) 上恰有一个简单谱点，且其余非零谱满足 \(|\mu|\neq \sqrt{\lambda}\)，则 \(\mathcal G_M\) 仅在 \(\sigma=\frac12\) 处出现唯一折角。
\end{theorem}
\begin{proof}
由命题 \ref{prop:finite-rh-phase-lift-artin}，对任意 \(z\in\CC\) 有
\[
\det(I-zM^{\langle q\rangle})
=
\prod_{k=0}^{q-1}\det(I-z\omega_q^k M).
\]
取
\[
z=\lambda^{-\sigma}e^{i\theta},
\]
并对两边取绝对值对数、再在 \(\theta\in[0,2\pi]\) 上积分，可得
\[
\mathcal G_{M^{\langle q\rangle}}(\sigma)
=
\frac1q\sum_{k=0}^{q-1}
\int_0^{2\pi}
\log\left|\det\!\left(I-\lambda^{-\sigma}e^{i(\theta+2\pi k/q)}M\right)\right|
\frac{d\theta}{2\pi}.
\]
由于积分区间是完整圆周，变量平移 \(\theta\mapsto \theta+2\pi k/q\) 不改变平均值，故上式等于
\[
\int_0^{2\pi}
\log\left|\det\!\left(I-\lambda^{-\sigma}e^{i\theta}M\right)\right|
\frac{d\theta}{2\pi}.
\]
这就证明了循环相位细分在归一化后完全不可见。

再把行列式按非零特征值分解为
\[
\det\!\left(I-\lambda^{-\sigma}e^{i\theta}M\right)
=
\prod_{j\in J}\left(1-\lambda^{-\sigma}e^{i\theta}\mu_j\right),
\]
并逐项应用 Jensen 公式
\[
\int_0^{2\pi}\log|1-re^{i\theta}|\frac{d\theta}{2\pi}=\log^+ r
\qquad (r\ge 0),
\]
即可得到
\[
\int_0^{2\pi}
\log\left|\det\!\left(I-\lambda^{-\sigma}e^{i\theta}M\right)\right|
\frac{d\theta}{2\pi}
=
\sum_{j\in J}\log^+\!\left(|\mu_j|\lambda^{-\sigma}\right),
\]
亦即题述分段线性公式。

最后，每一项
\[
\max\!\bigl(0,\log|\mu_j|-\sigma\log\lambda\bigr)
\]
的折点位于
\[
\sigma=\frac{\log|\mu_j|}{\log\lambda}.
\]
若仅有一个简单谱点满足 \(|\mu_j|=\sqrt{\lambda}\)，则唯一落在 \(\sigma=\frac12\) 的折点就是该项，其余项的折点全部远离 \(\frac12\)。故 \(\mathcal G_M\) 在 \(\sigma=\frac12\) 处仅有唯一折角。证毕。
\end{proof}

\begin{theorem}[有限窗梳齿拟合与全局去晶格化的有理无关障碍]\label{thm:conclusion-cyclic-lift-finitewindow-delattice-obstruction}
设 \(M_1,\dots,M_A\) 为有限个 RH-sharp 核，Perron 根分别为 \(\lambda_a>1\)。对每个 \(a\)，取一条循环 lift 阶 \(q_a\ge 2\) 与一个达界模态
\[
\mu_a=\sqrt{\lambda_a}\,e^{i\theta_a},
\]
并记其在临界线 \(\Re(s)=\frac12\) 上诱导的虚部梳齿为
\[
\mathcal C_a:=
\left\{
\frac{\theta_a+2\pi n}{q_a\log\lambda_a}
\;:\;
n\in\ZZ
\right\},
\qquad
\Delta_a:=\frac{2\pi}{q_a\log\lambda_a}.
\]
则有：
\begin{enumerate}
  \item 对任意单个 RH-sharp 核 \(M\)、任意有限集合 \(\Gamma\subset\RR\) 与任意 \(\varepsilon>0\)，存在某个循环阶 \(q\) 使得对每个 \(\gamma\in\Gamma\)，都存在 \(\tau\in\mathcal C_q\) 满足
  \[
  |\gamma-\tau|<\varepsilon.
  \]
  换言之，单尺度循环提升已经足以实现任意有限窗内的临界线梳齿逼近。
  \item 若对所有 \(a,b\) 都有
  \[
  \frac{\log\lambda_a}{\log\lambda_b}\in\QQ,
  \]
  则存在 \(T>0\) 使得
  \[
  \bigcup_{a=1}^{A}\mathcal C_a
  \subseteq
  \bigcup_{a=1}^{A}\bigl(\alpha_a+T\ZZ\bigr)
  \]
  对某些实数 \(\alpha_a\) 成立。特别地，全部临界线极点仍被压在有限条等差竖列上，故保留一个非零共同虚周期。
  \item 反之，若存在某对 \(a,b\) 满足
  \[
  \frac{\log\lambda_a}{\log\lambda_b}\notin\QQ,
  \]
  则不存在任何 \(T>0\) 使 \(\bigcup_{a=1}^{A}\mathcal C_a\) 可写成有限个 \(T\ZZ\) 的平移并。因而有限通道情形下的全局去晶格化至少需要对数 Perron 根的有理无关性。
\end{enumerate}
\end{theorem}
\begin{proof}
第一条由推论 \ref{cor:zeta-cyclic-lift-dirichlet-comb} 的梳齿公式直接给出。对单个 RH-sharp 核，临界线上的虚部网格间距为
\[
\Delta_q=\frac{2\pi}{q\log\lambda}.
\]
取 \(q\) 充分大使 \(\Delta_q<2\varepsilon\)。则对任意 \(\gamma\in\Gamma\)，选取最接近 \(\gamma\) 的格点
\[
\tau=\frac{\theta+2\pi n}{q\log\lambda}\in\mathcal C_q
\]
即可得到
\[
|\gamma-\tau|\le \frac{\Delta_q}{2}<\varepsilon.
\]

对第二条，若 \(\log\lambda_a/\log\lambda_b\in\QQ\) 对所有 \(a,b\) 成立，则
\[
\frac{\Delta_a}{\Delta_1}
=
\frac{q_1}{q_a}\cdot\frac{\log\lambda_1}{\log\lambda_a}
\in\QQ
\qquad (1\le a\le A).
\]
把这些有理数写成既约分数 \(m_a/n_a\)，并令
\[
L:=\mathrm{lcm}(n_1,\dots,n_A),
\qquad
T:=\frac{\Delta_1}{L}.
\]
于是对每个 \(a\)，
\[
\Delta_a=\frac{m_aL}{n_a}\,T\in T\NN.
\]
因此每一条梳齿
\[
\mathcal C_a=\alpha_a+\Delta_a\ZZ
\]
都包含于某条 \(T\)-晶格平移 \(\alpha_a+T\ZZ\) 中，从而整个并集仍是有限条 \(T\)-等差竖列的并。

对第三条，反设存在 \(T>0\) 与有限多平移 \(\beta_1+T\ZZ,\dots,\beta_r+T\ZZ\)，使
\[
\bigcup_{a=1}^{A}\mathcal C_a\subseteq \bigcup_{j=1}^{r}(\beta_j+T\ZZ).
\]
固定某个 \(a\)。由于 \(\mathcal C_a\) 是无限集而右端只有有限条 \(T\)-晶格，鸽巢原理保证存在某个 \(j\) 使
\[
\mathcal C_a\cap(\beta_j+T\ZZ)
\]
为无限集。取其中两点
\[
\alpha_a+n\Delta_a,\qquad \alpha_a+n'\Delta_a
\]
落在同一平移类内，则其差满足
\[
(n-n')\Delta_a\in T\ZZ.
\]
故 \(\Delta_a/T\in\QQ\)。同理 \(\Delta_b/T\in\QQ\)，于是
\[
\frac{\Delta_a}{\Delta_b}\in\QQ.
\]
但
\[
\frac{\Delta_a}{\Delta_b}
=
\frac{q_b}{q_a}\cdot\frac{\log\lambda_b}{\log\lambda_a},
\]
从而推出 \(\log\lambda_a/\log\lambda_b\in\QQ\)，与假设矛盾。故不存在这样的共同晶格周期 \(T\)。证毕。
\end{proof}

\begin{corollary}[单临界相位类的单梳齿与奇支撑平方根异常]\label{cor:conclusion-realinput40-single-comb-odd-support-chain}
对真实输入 \(40\) 状态核，隐藏的 \((-A_0)\) 因子所强制的唯一达界模态
\[
\mu_\ast=-\sqrt{\lambda}=-\varphi
\]
同时诱导以下三种可见现象：
\begin{enumerate}
  \item 临界线上的非平凡极点只形成一条半格点梳齿；
  \item primitive 轨道计数中的 \(\lambda^{n/2}\) 级异常完全支撑在奇长度上；
  \item 其余全部非 Perron 模态都严格亚临界，因而不再贡献第二条临界梳齿。
\end{enumerate}
\end{corollary}
\begin{proof}
由定理 \ref{thm:conclusion-realinput40-single-channel-comb-rigidity}，真实输入 \(40\) 状态核除 Perron 根外恰有唯一达界模态 \(-\sqrt{\lambda}\)，且该模态完全来自隐藏的 \((-A_0)\) 因子。再由推论 \ref{cor:zeta-signflip-half-lattice}，这一符号翻转模态恰对应于临界线上的半格点梳齿。最后，定理 \ref{thm:conclusion-realinput40-primitive-sqrt-anomaly-odd-support} 表明同一模态在 primitive 层恰产生
\[
-\mathbf 1_{\{n\ \mathrm{odd}\}}\frac{\lambda^{n/2}}{n}
\]
这一奇支撑平方根异常，而偶长度上的同尺度项被 Möbius 反演中的 \(d=2\) 项精确抵消。故三种现象由同一单临界相位类同时锁定。证毕。
\end{proof}

\begin{theorem}[正性 Fredholm 完成中的边界模态簇不灭性]\label{thm:conclusion-positive-fredholm-boundary-obstruction}
固定 \(\lambda>1\)。设
\[
\zeta_{\mathcal D}(r)=\prod_{j\in J}(1-\beta_j r^{n_j})^{-m_j}
\]
为一组原子 Euler 数据所定义的 Euler 型乘积，并设存在迹类算子 \(L\) 与 \(r_0>0\)，使得
\[
\det(I-rL)^{-1}=\zeta_{\mathcal D}(r)
\qquad (|r|<r_0).
\]
若某个原子满足
\[
|\beta_j|^{1/n_j}=\sqrt{\lambda},
\]
则对任意 \(\alpha_j\in\CC\) 满足 \(\alpha_j^{n_j}=\beta_j\)，根集合
\[
\{\alpha_j\omega:\ \omega^{n_j}=1\}
\]
中的每一点都必须作为 \(L\) 的非零特征值出现，且代数重数至少为 \(m_j\)。

特别地，若该 Fredholm 实现同时满足注 \ref{rem:positive-fredholm-htf} 的正性/整数重数规格，则这一边界模态簇不可能通过同点主部消去而消失。于是相对于同一 Euler 接口，strict 几何不可能在正性 Fredholm 完成之后才被获得；若 strict 化发生，它必须发生在原子 Euler 接口被固定之前。
\end{theorem}
\begin{proof}
命题 \ref{prop:cyclic-euler-spectral-rigidity} 已经给出：一旦
\[
\det(I-rL)^{-1}=\zeta_{\mathcal D}(r),
\]
则 \(\zeta_{\mathcal D}\) 的零点结构强制 \(L\) 的非零谱多重集。对每个原子 \((\beta_j,n_j,m_j)\)，任选 \(\alpha_j^{n_j}=\beta_j\)，则集合
\[
\{\alpha_j\omega:\ \omega^{n_j}=1\}
\]
中的每个点都必须作为 \(L\) 的非零特征值出现，且代数重数至少为 \(m_j\)。因此若某个原子位于边界模长 \(\sqrt{\lambda}\)，则这一整簇边界模态已经被零点结构刚性锁定；若不同原子在同一点相遇，其代数重数只会相加，而不会相互抵消。

另一方面，注 \ref{rem:positive-fredholm-htf} 的正性 Fredholm-HTF 规格要求离散谱点具有整数代数重数 \(m_\lambda=\Tr(P_\lambda)\in\NN\)，并且相关主部系数由 \(\Tr(P_\lambda)\) 或更一般的 \(\Tr(WP_\lambda)\ge 0\) 给出。故在该规格下，同点主部消去被结构性排除。于是，只要边界原子已经进入 Euler 接口，任何保持该接口的正性 Fredholm 完成都不能将其抹除；因此 strict 化若要发生，只能发生在 Euler 接口被固定之前。证毕。
\end{proof}

\begin{theorem}[universal solenoid 的有限见证与寄存器谱刚性]\label{thm:conclusion-universal-solenoid-finite-witness-spectral-rigidity}
设 \(\Sigma_{\QQ}\) 为通用 solenoid，\(\kappa\) 为其上由有限个有理频率特征标生成的三角多项式。并假设 \(\kappa\) 的寄存器侧 Euler 接口由一组只使用 \(\kappa\) 所涉分母素因子的复权原子数据
\[
\mathcal D=\{(\beta_j,n_j,m_j)\}_{j\in J}\subset \CC\times\NN_{\ge 1}\times\NN
\]
给出，且满足
\[
\sum_{j\in J}m_jn_j|\beta_j|^{1/n_j}<\infty,
\qquad
|\beta_j|<1.
\]
记
\[
\zeta_{\mathcal D}(r):=\prod_{j\in J}(1-\beta_j r^{n_j})^{-m_j}.
\]
则存在有限素数集 \(S\) 与 \(\Sigma_S\) 上的三角多项式 \(\kappa_S\)，使得
\[
\kappa=\kappa_S\circ\pi_S.
\]
并且同一 \(S\) 已固定上述寄存器侧 Euler 数据；由此可构造显式迹类算子 \(L_{\mathcal D}\) 使
\[
\det(I-rL_{\mathcal D})^{-1}=\zeta_{\mathcal D}(r)
\qquad (|r|<1).
\]
对任意其他满足同一 Fredholm \(\zeta\)-接口的迹类实现 \(L\)，其非零谱多重集与 \(L_{\mathcal D}\) 完全一致。因而一切由寄存器侧原子谱支配的边界模态、临界梳齿与 Toeplitz/Gram/正定性证书，都已在某个有限 prime-register 阶段被完全见证；从 \(\Sigma_S\) 到 \(\Sigma_{\QQ}\) 的极限过程不会创造新的寄存器侧谱原子。另一方面，相位--Witt 转移算子在 Mellin 频域仅以
\[
\mathcal M(\mathcal T G)(s)=\zeta(s)\,\mathcal M G(s)
\]
这一公共乘子作用，故 universal 层至多叠加与 \(S\) 无关的相位乘子，而不改变寄存器侧原子支撑。
\end{theorem}
\begin{proof}
由定理 \ref{thm:cdim-qhat-finite-prime-witness}，存在有限素数集 \(S\) 与 \(\Sigma_S\) 上的三角多项式 \(\kappa_S\) 使
\[
\kappa=\kappa_S\circ\pi_S,
\]
并且凡由这一有限频率族生成的 Toeplitz、Gram 与正定性证书，也都经由同一 \(\pi_S\) 因子化。由于定理假设寄存器侧 Euler 接口只使用这些分母素因子，故同一 \(S\) 已经固定原子数据 \(\mathcal D\)。

接着，由推论 \ref{cor:cyclic-euler-product-complex-section}，\(\mathcal D\) 具有显式循环块迹类实现 \(L_{\mathcal D}\)，满足
\[
\det(I-rL_{\mathcal D})^{-1}=\zeta_{\mathcal D}(r)
\qquad (|r|<1).
\]
再由命题 \ref{prop:cyclic-euler-spectral-rigidity}，任何其他满足同一 Fredholm \(\zeta\)-接口的迹类实现 \(L\) 都拥有与 \(L_{\mathcal D}\) 完全相同的非零谱多重集。因而凡由寄存器侧原子谱决定的边界模态与临界梳齿，都已经在有限阶段 \(\Sigma_S\) 被完全固定。

最后，由相位--Witt 转移算子的 Mellin 对角化公式 \eqref{eq:scale-mellin-diagonalization}，
\[
\mathcal M(\mathcal T G)(s)=\zeta(s)\,\mathcal M G(s),
\]
可见从相位侧引入的 universal 极限只是在频域乘上一个与 \(S\) 无关的公共因子 \(\zeta(s)\)。因此，极限过程不会生成新的寄存器侧原子支撑；它至多在相位侧叠加一个公共 Mellin 乘子。证毕。
\end{proof}

\begin{corollary}[strict 几何的共尾有限性]\label{cor:conclusion-strict-geometry-cofinal-finite}
在定理 \ref{thm:conclusion-positive-fredholm-boundary-obstruction} 与定理 \ref{thm:conclusion-universal-solenoid-finite-witness-spectral-rigidity} 的共同设定下，固定 \(\lambda>1\)，并称原子 \((\beta,n,m)\) 为边界寄存器原子，若
\[
|\beta|^{1/n}=\sqrt{\lambda}.
\]
对任意有限生成观测 \(\kappa\)，下列两命题等价：
\[
\text{(i) 完成到 }\Sigma_{\QQ}\text{ 后不含任何边界寄存器原子;}
\]
\[
\text{(ii) 在某个有限见证阶段 }\Sigma_S\text{ 上已不含任何边界寄存器原子。}
\]
若某个有限见证阶段已经出现边界寄存器原子，则任何保持同一 Euler 接口的正性 Fredholm 完成都不能将其抹除。因而对有限生成观测而言，strict 临界几何是一个共尾有限性质，而 RH-sharp 型失效则是有限可见且在完成层不可修复的。
\end{corollary}
\begin{proof}
若某个有限见证阶段 \(\Sigma_S\) 上已经出现边界寄存器原子，则由定理 \ref{thm:conclusion-universal-solenoid-finite-witness-spectral-rigidity}，其对应的寄存器侧谱原子在通用层不会消失；故 (i) 蕴含 (ii)。

反过来，若 \(\Sigma_S\) 上不含任何边界寄存器原子，则同一定理说明从 \(\Sigma_S\) 过渡到 \(\Sigma_{\QQ}\) 的极限过程不会创造新的寄存器侧谱原子，故 (ii) 蕴含 (i)。

最后一段直接应用定理 \ref{thm:conclusion-positive-fredholm-boundary-obstruction}：一旦某个有限见证阶段已经固定了边界原子，任何保持同一 Euler 接口的正性 Fredholm 完成都不能将其抹除。证毕。
\end{proof}

\begin{theorem}[经典 RH 的 arithmetic solenoid 重述]\label{thm:conclusion-classical-rh-arithmetic-solenoid-restatement}
记一维算术 solenoid
\[
\Sigma_{\mathrm{sol}}
:=
(\RR\times \widehat{\ZZ})/\ZZ
\cong
\mathbb A_{\QQ}/\QQ
\cong
\widehat{\QQ},
\]
并记其相位圆商为
\[
\pi_{\mathrm{ph}}:\Sigma_{\mathrm{sol}}\twoheadrightarrow \TT.
\]
沿用式 \eqref{eq:app-s-of-w} 的圆盘截面
\[
s(w)=\frac12+\mathrm{i}\,\mathcal{C}^{-1}(w)
\qquad (|w|<1),
\]
定义 solenoid 相位边界上的完成截面
\[
Z_{\Sigma}(w):=\Xi_{\zeta}\!\bigl(s(w)\bigr),
\]
以及其规范 Carath\'eodory 对数导数
\[
\mathscr C_{\Sigma}(w)
:=
-\frac{\Xi_\zeta'(s(w))}{\Xi_\zeta(s(w))}
=
\mathscr C_\zeta(w).
\]
则下列陈述等价：
\begin{enumerate}
  \item 经典 RH 成立；
  \item 命题 \(\mathbf{RH}_{\Sigma}\) 成立：非平凡零点作为 \(\Sigma_{\mathrm{sol}}\) 上完成 \(\Xi\)-截面的谱支撑，全部集中于唯一的幺正 Mellin 切片
  \[
  \Re(s)=\frac12;
  \]
  \item
  \[
  Z_{\Sigma}(w)\neq 0
  \qquad (|w|<1);
  \]
  \item \(\mathscr C_{\Sigma}\) 属于 Carath\'eodory 类；
  \item 若
  \[
  \mathscr C_{\Sigma}(w)=c_0+2\sum_{n\ge 1}c_n w^n,
  \qquad c_{-n}:=\overline{c_n},
  \]
  则对一切 \(N\ge 0\)，Toeplitz 主子矩阵
  \[
  \mathbf T_N:=(c_{j-k})_{0\le j,k\le N}
  \]
  都满足
  \[
  \mathbf T_N\succeq 0.
  \]
\end{enumerate}
特别地，在这一语言中，prime-register 的 profinite 纤维 \(\widehat{\ZZ}\) 只承担寄存与延拓方向，而不生成任何偏离唯一幺正切片的额外解析谱支撑。故上述 \(\mathbf{RH}_{\Sigma}\) 只是经典 RH 的等价重写，而非现成证明。
\end{theorem}
\begin{proof}
由定理 \ref{thm:conclusion-finite-prime-solenoid-terminal-object}，有限素数局部化的 solenoid 已经表现出同一刚性特征：其全部连通 Lie 可见部分都压缩到唯一圆商 \(\TT\)，而 prime-register 信息则滞留在全不连通核中。推论 \ref{cor:conclusion-fibonacci-solenoid-dual-q} 进一步把通用一维 solenoid 识别为
\[
\widehat{\QQ}\cong \mathbb A_{\QQ}/\QQ,
\]
故 \(\Sigma_{\mathrm{sol}}\) 正是这一 arithmetic universal object 的标准模型。于是其解析可见方向只有相位圆商 \(\pi_{\mathrm{ph}}\)，而 profinite 纤维只提供寄存器侧的外置方向。

另一方面，由相位--Witt 转移算子的 Mellin 对角化公式 \eqref{eq:scale-mellin-diagonalization}，
\[
\mathcal M(\mathcal T G)(s)=\zeta(s)\,\mathcal M G(s),
\]
可见完成化后的解析读出沿 Mellin 变量只引入单一复参数 \(s\)；因此“全部非平凡零点集中于 \(\Re(s)=\tfrac12\)”正是该一维相位读出上的 unitary-slice 断言。若存在偏离临界线的零点，则其在圆盘完成化下恰对应于 \(|w|<1\) 内的内点零点；反之，盘内零点又唯一回升为偏离 \(\Re(s)=\tfrac12\) 的解析谱支撑。

具体地，定理 \ref{thm:app-rh-iff-disk-zero-free} 已证明
\[
\mathrm{RH}
\quad\Longleftrightarrow\quad
Z_\Sigma(w)=\Xi_\zeta(s(w))\neq 0
\qquad (|w|<1),
\]
这给出 (1)\(\Leftrightarrow\)(3)。而推论 \ref{cor:app-horizon-carath-rh} 说明上式又等价于
\[
\mathscr C_\Sigma=\mathscr C_\zeta
\]
属于 Carath\'eodory 类，故 (3)\(\Leftrightarrow\)(4)。再由定理 \ref{thm:app-horizon-toeplitz-lmi}，
\[
\mathscr C_\Sigma\ \text{为 Carath\'eodory}
\quad\Longleftrightarrow\quad
\mathbf T_N\succeq 0\ \ (\forall N\ge 0),
\]
于是 (4)\(\Leftrightarrow\)(5)。

最后，由上半平面到圆盘的 Cayley 完成化把偏离临界线的零点与盘内内点零点一一对应，而 \(\Sigma_{\mathrm{sol}}\) 的唯一连通解析读出层正是相位圆所给出的 Mellin 切片，故 (2) 与 (3) 只是同一事实的 solenoid/圆盘两种表述。综上，(1)--(5) 全部等价，且最后一段只是对这一等价链的结构性读法。证毕。
\end{proof}

\begin{theorem}[\(\varphi\)-adic Mellin 指纹的极点格刚性]\label{thm:conclusion-phiadic-mellin-pole-lattice-rigidity}
设 \(\mathcal M(s)\) 为主稿第 \(13.145\) 节中的 Mellin 指纹。若其在变量
\[
z=\varphi^{-s}
\]
上具有有理表示
\[
\mathcal M(s)=R(\varphi^{-s}),
\]
其中 \(R(z)\) 为有理函数，则：
\begin{enumerate}
  \item \(\mathcal M\) 具有纯虚周期
  \[
  \mathcal M\!\left(s+\frac{2\pi i}{\log\varphi}\right)=\mathcal M(s);
  \]
  \item 任一极点 \(s_0\) 强制整条垂直格点轨道
  \[
  s_0+\frac{2\pi i}{\log\varphi}\ZZ
  \]
  全部为极点，且主部完全相同；
  \item 因而有限状态 \(\varphi\)-adic 编译不允许出现孤立的非格点极点。
\end{enumerate}
\end{theorem}
\begin{proof}
由
\[
\varphi^{-(s+\frac{2\pi i}{\log\varphi})}
=
\varphi^{-s}e^{-2\pi i}
=
\varphi^{-s},
\]
立即得到
\[
\mathcal M\!\left(s+\frac{2\pi i}{\log\varphi}\right)
=
R\!\left(\varphi^{-(s+\frac{2\pi i}{\log\varphi})}\right)
=
R(\varphi^{-s})
=
\mathcal M(s),
\]
即第一条。

若 \(s_0\) 是 \(\mathcal M\) 的极点，则 \(z_0:=\varphi^{-s_0}\) 是 \(R\) 的极点。对任意整数 \(k\)，都有
\[
\varphi^{-(s_0+\frac{2\pi i}{\log\varphi}k)}=z_0.
\]
因此
\[
\mathcal M\!\left(s_0+\frac{2\pi i}{\log\varphi}k\right)=R(z_0)
\]
落在同一个有理函数极点上，故整条轨道全部为极点，且其 Laurent 主部由同一 \(R\) 在 \(z_0\) 处的极点数据给出，完全一致。

若存在孤立的非格点极点，则其纯虚平移轨道应只含有限个点，这与上一步推出的整条垂直格点轨道矛盾。故有限状态 \(\varphi\)-adic 编译下不可能出现孤立非格点极点。证毕。
\end{proof}
\leanverified{paper\_conclusion\_phiadic\_mellin\_pole\_lattice\_rigidity}

\begin{corollary}[\(\varphi\)-adic 有限编译的孤立极点禁阻]\label{cor:conclusion-phiadic-finite-compile-no-isolated-pole}
在定理 \ref{thm:conclusion-phiadic-mellin-pole-lattice-rigidity} 的设定下，任何试图用有限维 \(\varphi\)-adic 矩阵细分或有限状态编译来制造单个异常极点的方案都必失败。更准确地说，若某 Mellin 指纹在某条垂直带内存在唯一孤立极点 \(s_0\)，则它不可能来自主稿第 \(13.145\) 节所允许的有限编译机制。
\end{corollary}
\begin{proof}
若该 Mellin 指纹来自主稿允许的有限编译机制，则按定理 \ref{thm:conclusion-phiadic-mellin-pole-lattice-rigidity}，极点 \(s_0\) 必强制整条轨道
\[
s_0+\frac{2\pi i}{\log\varphi}\ZZ
\]
全部为极点。于是 \(s_0\) 不可能是某条垂直带内的唯一孤立极点，矛盾。故此类单极点机制不可能由有限 \(\varphi\)-adic 编译产生。证毕。
\end{proof}

\begin{theorem}[共尾固定商不可能性]\label{thm:conclusion-artin-visible-hull-no-cofinal-fixed-quotient}
设 \(\{D_m\}_{m\ge 1}\) 是一列分辨率递增的数据，每一级 \(D_m\) 都定义在有限群 \(G_m\) 上，并记其 Artin-visible hull 为
\[
Q_m:=Q_\ast(D_m).
\]
若存在某个固定有限群 \(H\) 以及一列满同态
\[
\pi_m:G_m\twoheadrightarrow H
\]
使得对所有充分大的 \(m\)，\(\pi_m\) 都保持 \(D_m\) 的全部 determinant/类函数信息，则每个充分大的 \(m\) 都存在满同态
\[
H\twoheadrightarrow Q_m.
\]
因此必有
\[
|Q_m|\le |H|
\qquad (m\gg 1),
\]
并且每个 \(Q_m\) 的全部合成因子都必须来自 \(H\)。

特别地，只要下列任一情形发生，就不存在这样的固定有限商 \(H\)：
\begin{enumerate}
  \item \(|Q_m|\to\infty\)；
  \item \(Q_m\) 的合成因子最终逃离任意有限清单；
  \item \(Q_m\) 非最终阿贝尔，而 \(H\) 取阿贝尔。
\end{enumerate}
\leanverified{paper\_conclusion\_cofinal\_fixed\_quotient\_seeds}
\leanverified{paper\_conclusion\_cofinal\_fixed\_quotient\_package}
\end{theorem}

\begin{proof}
由定理 \ref{thm:conclusion-artin-determinant-classfunction-joint-minimal-quotient}，\(Q_m\) 是保持第 \(m\) 级全部活跃信息的唯一最小商。
故若 \(\pi_m\) 保持第 \(m\) 级信息，则按唯一最小性，\(\pi_m\) 必经由 \(Q_m\) 因子化，即存在满同态
\[
H\twoheadrightarrow Q_m.
\]
于是
\[
|Q_m|\le |H|.
\]
而商群的合成因子皆出自母群，故 \(Q_m\) 的合成因子必来自 \(H\)。
三条否定条件随即成立。证毕。
\end{proof}

\begin{conjecture}[共尾 RH 证书的非阿贝尔可见壳增长]\label{conj:conclusion-solenoidal-rh-visible-hull-growth}
在 solenoidal--Mellin--Fredholm 语义下，设 \(\mathcal K_m\) 为第 \(m\) 级群标记核，\(\mathcal E_m\) 为维持终端证书链所必需的有限类函数族，至少包含主导二主项见证函数 \(\psi_{\ast,u,m}\) 的适当有限抽样。
记相应的 Artin-visible hull 为
\[
Q_m:=Q_\ast(\mathcal K_m,\mathcal E_m).
\]
则应有：
\begin{enumerate}
  \item \(\{Q_m\}\) 不应最终有界；
  \item \(\{Q_m\}\) 不应最终阿贝尔；
  \item 更强地，\(\{Q_m\}\) 的非阿贝尔部分应在共尾分辨率上持续产生新的合成因子。
\end{enumerate}
\end{conjecture}

这一命题把 screen--Artin 双障碍压缩为单一群论断言：任何共尾完备的 RH 终端证书，都必须携带一个真正增长且真正非阿贝尔的可见壳；不存在固定有限商可以长期承载全部活跃 determinant 因子与全部主导见证函数。

\endinput

\begin{theorem}[有理主弧 Gram 核的 Ramanujan--Dirichlet 双包恒等式]\label{thm:conclusion-majorarc-gram-ramanujan-dirichlet-doublepacket}
\leanverified{paper\_conclusion\_majorarc\_gram\_ramanujan\_dirichlet\_doublepacket}
对整数 \(M\ge 2\) 与 \(L\asymp M^2\)，记
\[
Q_M:=\left\{\frac{a}{q}:1\le q\le M,\ (a,q)=1\right\}\subset \QQ/\ZZ,
\]
并令 \(K_{M,L}\) 为附录 B.6.1 中的有理主弧 Gram 核算子。则有如下两条精确层析恒等式。

若 \(g_\alpha\in \ell^2(Q_M)\) 在每个分母层上为常数，
\[
g_\alpha(a/q)=\alpha_q,
\]
则
\[
\langle g_\alpha,K_{M,L}g_\alpha\rangle
=
\sum_{n\in\ZZ}\psi^{\flat}(n/L)
\left|
\sum_{q\le M}\alpha_q\,c_q(n)
\right|^2.
\]

若对每个 \(q\le M\) 选定一个本原 Dirichlet 角色 \(\chi_q\pmod q\)，并令
\[
g_{\beta,\chi}(a/q):=\beta_q\chi_q(a),
\]
则
\[
\langle g_{\beta,\chi},K_{M,L}g_{\beta,\chi}\rangle
=
\sum_{n\in\ZZ}\psi^{\flat}(n/L)
\left|
\sum_{q\le M}\beta_q\,\tau(\chi_q)\chi_q(n)\mathbf 1_{(n,q)=1}
\right|^2.
\]

故 Gram 核在“分母层常值”子空间上精确化为 Ramanujan 包，而在“分母层单角色”子空间上精确化为 Dirichlet 波包。
\end{theorem}
\begin{proof}
把 \(g=g_\alpha\) 代入附录 B.6.1 的 Gram 恒等式，得到
\[
\langle g_\alpha,K_{M,L}g_\alpha\rangle
=
\sum_{n\in\ZZ}\psi^{\flat}(n/L)
\left|
\sum_{q\le M}\alpha_q
\sum_{a\in U(q)}e^{2\pi ian/q}
\right|^2.
\]
而内层和正是 Ramanujan 和 \(c_q(n)\)，故得第一式。

同理，对 \(g_{\beta,\chi}\) 有
\[
\langle g_{\beta,\chi},K_{M,L}g_{\beta,\chi}\rangle
=
\sum_{n\in\ZZ}\psi^{\flat}(n/L)
\left|
\sum_{q\le M}\beta_q
\sum_{a\in U(q)}\chi_q(a)e^{2\pi ian/q}
\right|^2.
\]
再用正文中 primitive Gauss--Dirichlet 展开
\[
\sum_{a\in U(q)}\chi_q(a)e^{2\pi ian/q}
=
\tau(\chi_q)\chi_q(n)\mathbf 1_{(n,q)=1},
\]
即得第二式。证毕。
\end{proof}

\begin{theorem}[单尺度谱降驱动的全尺度 Ramanujan--Dirichlet 波包压制]\label{thm:conclusion-majorarc-single-scale-drop-forces-ramanujan-dirichlet-suppression}
若附录 B.6.1 的单尺度谱降假设成立，即存在某个 \(M_0\) 与 \(L_0\asymp M_0^2\) 使
\[
\|K_{M_0,L_0}\|_{\mathrm{op}}\le M_0^{-\eta}
\qquad (\eta>0),
\]
则存在 \(\eta'>0\)，使对一切 \(M\ge M_0\) 与同型标度 \(L\asymp M^2\)，有
\[
\sum_{n\in\ZZ}\psi^{\flat}(n/L)
\left|
\sum_{q\le M}\alpha_q\,c_q(n)
\right|^2
\ll
M^{-\eta'}
\sum_{q\le M}\varphi(q)|\alpha_q|^2,
\]
以及
\[
\sum_{n\in\ZZ}\psi^{\flat}(n/L)
\left|
\sum_{q\le M}\beta_q\,\tau(\chi_q)\chi_q(n)\mathbf 1_{(n,q)=1}
\right|^2
\ll
M^{-\eta'}
\sum_{q\le M}\varphi(q)|\beta_q|^2.
\]

特别地，若取 \(\alpha_q=C_q(N)\) 为某个 primitive 核 \(N\) 的 cyclotomic 原子，则
\[
\sum_{n\in\ZZ}\psi^{\flat}(n/L)
\left|
\sum_{q\le M}C_q(N)c_q(n)
\right|^2
\ll
M^{-\eta'}
\sum_{q\le M}\varphi(q)|C_q(N)|^2.
\]
\leanverified{paper\_conclusion\_majorarc\_single\_scale\_drop\_forces\_ramanujan\_dirichlet\_suppression}
\end{theorem}
\begin{proof}
由定理 \ref{thm:conclusion-majorarc-gram-ramanujan-dirichlet-doublepacket}，
上述左端都可写为
\[
\langle g,K_{M,L}g\rangle
\le
\|K_{M,L}\|_{\mathrm{op}}\|g\|_2^2.
\]
而对两类波包都成立
\[
\|g\|_2^2
=
\sum_{q\le M}\sum_{a\in U(q)}|{\rm coefficient}_q|^2
=
\sum_{q\le M}\varphi(q)|{\rm coefficient}_q|^2.
\]
再由附录 B.6.1 的尺度自举与单尺度谱降传播，
\[
\|K_{M,L}\|_{\mathrm{op}}\ll M^{-\eta'},
\]
代回即得两条全尺度压制式。取 \(\alpha_q=C_q(N)\) 便得到最后一式。证毕。
\end{proof}

\begin{proposition}[有限审计足以触发全尺度 Dirichlet 波包压制]\label{prop:conclusion-majorarc-finite-audit-suffices-for-global-wavepacket-suppression}
\leanverified{paper\_conclusion\_majorarc\_finite\_audit\_suffices\_for\_global\_wavepacket\_suppression}
对固定 primitive sofic--权系统，附录 B.6.1 给出显式可计算的有限审计集
\[
\mathcal T\subset (0,2\pi),
\qquad
\mathcal M\subset \NN,
\]
使统一扭曲谱隙可归约为有限多对 \((\theta,M)\) 的核验。因而，在附录 B.6.1 的制度化框架内，定理 \ref{thm:conclusion-majorarc-single-scale-drop-forces-ramanujan-dirichlet-suppression} 的全尺度 Ramanujan--Dirichlet 波包压制，可由有限审计程序触发。
\end{proposition}
\begin{proof}
附录 B.6.1 的有限审计定理断言：统一扭曲谱隙成立，当且仅当有限集合 \(\mathcal T\times \mathcal M\) 上的相应谱核验全部通过。另一方面，附录 B.6.1 的单尺度谱降传播定理又说明，一旦得到这种统一谱降，它便沿 \(L\asymp M^2\) 的尺度族传播到全部 \(M\)。两者合并，即得全尺度 Ramanujan--Dirichlet 波包压制并不需要无穷多尺度逐一验证，而只需有限审计。证毕。
\end{proof}

\begin{theorem}[素数导子中心化同余切片的 Gram--Fourier 精确公式]\label{thm:conclusion-prime-conductor-centered-slice-gram-fourier-identity}
\leanverified{paper\_conclusion\_prime\_conductor\_centered\_slice\_gram\_fourier\_identity}
设 \(q\ge 3\) 为素数，\(N\) 为 primitive 核。记
\[
A_N(\chi):=
\begin{cases}
L_N(\chi;q)/\tau(\chi), & \chi\neq \chi_0,\\[4pt]
\log C(N)-B_q(N), & \chi=\chi_0,
\end{cases}
\]
并定义单位同余类切片势能
\[
V_{q,r}(N):=\sum_{m\ge 0}\frac{u_{r+mq}(N)}{r+mq},
\qquad r\in U(q),
\]
及其质心
\[
\bar V_q(N):=\frac{1}{q-1}\sum_{r\in U(q)}V_{q,r}(N)
=
\frac{A_N(\chi_0)}{q-1}.
\]
再定义中心化角色包
\[
g_{N,q}^{\mathrm{cen}}(a/q):=
\sum_{\chi\neq \chi_0}\frac{A_N(\chi)}{\tau(\chi)}\chi(a),
\qquad a\in U(q),
\]
并在其余分母层取零。对任意 \(M\ge q\) 与 \(L>0\)，设
\[
W_{q,L}(r):=\sum_{\ell\in\ZZ}\psi^{\flat}\!\left(\frac{r+\ell q}{L}\right),
\qquad r\in U(q).
\]
则成立精确恒等式
\[
\langle g_{N,q}^{\mathrm{cen}},K_{M,L}g_{N,q}^{\mathrm{cen}}\rangle
=
(q-1)^2
\sum_{r\in U(q)}W_{q,L}(r)\,
\bigl|V_{q,r}(N)-\bar V_q(N)\bigr|^2.
\]
\end{theorem}
\begin{proof}
由定理 \ref{thm:conclusion-majorarc-gram-ramanujan-dirichlet-doublepacket}，
\[
\langle g_{N,q}^{\mathrm{cen}},K_{M,L}g_{N,q}^{\mathrm{cen}}\rangle
=
\sum_{n\in\ZZ}\psi^{\flat}(n/L)\,|S_q(n)|^2,
\]
其中
\[
S_q(n):=
\sum_{\chi\neq \chi_0}\frac{A_N(\chi)}{\tau(\chi)}
\sum_{a\in U(q)}\chi(a)e^{2\pi ian/q}.
\]
对非平凡角色，primitive Gauss 和公式给出
\[
\sum_{a\in U(q)}\chi(a)e^{2\pi ian/q}
=
\tau(\chi)\chi(n)\mathbf 1_{q\nmid n},
\]
故
\[
S_q(n)
=
\mathbf 1_{q\nmid n}\sum_{\chi\neq \chi_0}A_N(\chi)\chi(n).
\]
若 \(n\equiv r\pmod q\) 且 \(r\in U(q)\)，则由角色正交反演，
\[
\sum_{\chi\in \widehat{U(q)}}A_N(\chi)\chi(r)=(q-1)V_{q,r}(N),
\]
从而
\[
\sum_{\chi\neq \chi_0}A_N(\chi)\chi(r)
=(q-1)V_{q,r}(N)-A_N(\chi_0)
=(q-1)\bigl(V_{q,r}(N)-\bar V_q(N)\bigr).
\]
把 \(n=r+\ell q\) 重新分组，即得所示公式。证毕。
\end{proof}

\begin{corollary}[单位同余类中心化位势的加权方差坍缩律]\label{cor:conclusion-prime-conductor-centered-slice-weighted-variance-collapse}
在定理 \ref{thm:conclusion-prime-conductor-centered-slice-gram-fourier-identity} 的设定下，成立上界
\[
\sum_{r\in U(q)}W_{q,L}(r)\,
\bigl|V_{q,r}(N)-\bar V_q(N)\bigr|^2
\le
\frac{\|K_{M,L}\|_{\mathrm{op}}}{q(q-1)}
\sum_{\chi\neq \chi_0}|A_N(\chi)|^2.
\]

若进一步记
\[
w_*(q,L):=\min_{r\in U(q)}W_{q,L}(r)>0,
\]
则有无权方差控制
\[
\sum_{r\in U(q)}\bigl|V_{q,r}(N)-\bar V_q(N)\bigr|^2
\le
\frac{\|K_{M,L}\|_{\mathrm{op}}}{q(q-1)w_*(q,L)}
\sum_{\chi\neq \chi_0}|A_N(\chi)|^2.
\]

特别地，一旦定理 \ref{thm:conclusion-majorarc-single-scale-drop-forces-ramanujan-dirichlet-suppression} 的全尺度谱降 regime 被触发，则固定素数 \(q\) 的单位同余类切片中心化方差被迫按幂次衰减。
\end{corollary}
\begin{proof}
由定理 \ref{thm:conclusion-prime-conductor-centered-slice-gram-fourier-identity}，
\[
(q-1)^2
\sum_{r\in U(q)}W_{q,L}(r)\,
\bigl|V_{q,r}(N)-\bar V_q(N)\bigr|^2
=
\langle g_{N,q}^{\mathrm{cen}},K_{M,L}g_{N,q}^{\mathrm{cen}}\rangle
\le
\|K_{M,L}\|_{\mathrm{op}}\,
\|g_{N,q}^{\mathrm{cen}}\|_2^2.
\]
再由角色正交性，
\[
\|g_{N,q}^{\mathrm{cen}}\|_2^2
=
\sum_{a\in U(q)}
\left|
\sum_{\chi\neq \chi_0}\frac{A_N(\chi)}{\tau(\chi)}\chi(a)
\right|^2
=
(q-1)\sum_{\chi\neq \chi_0}\frac{|A_N(\chi)|^2}{|\tau(\chi)|^2}.
\]
因为 \(q\) 为素数而所有非平凡角色都本原，故
\[
|\tau(\chi)|^2=q.
\]
代回即得加权方差上界。再用 \(W_{q,L}(r)\ge w_*(q,L)\) 即得无权方差控制。最后一条由定理 \ref{thm:conclusion-majorarc-single-scale-drop-forces-ramanujan-dirichlet-suppression} 代入 \(\|K_{M,L}\|_{\mathrm{op}}\ll M^{-\eta'}\) 立得。证毕。
\end{proof}

\begin{conjecture}[Dirichlet 切片非坍缩迫使有限审计反复失效]\label{conj:conclusion-dirichlet-slice-noncollapse-forces-repeated-finite-audit-failure}
对固定 primitive sofic--权系统与固定素数 \(q\)，若存在常数 \(c>0\) 与无穷尺度列 \(M_j\to\infty\)、\(L_j\asymp M_j^2\)，使得
\[
\sum_{r\in U(q)}W_{q,L_j}(r)\,
\bigl|V_{q,r}(N)-\bar V_q(N)\bigr|^2
\ge
c\sum_{\chi\neq \chi_0}|A_N(\chi)|^2
\qquad (j\to\infty),
\]
则附录 B.6.1 的有限审计集 \(\mathcal T\times \mathcal M\) 中，必有某些审计对在无穷多个尺度上反复失效。换言之，单位同余类切片的持续非坍缩不可能与“有限审计永久通过”共存。
\end{conjecture}

\begin{theorem}[视界零知识的单侧收缩与指数区分器互斥]\label{thm:conclusion-hzom-one-sided-contraction-exponential-distinguisher-incompatibility}
\leanverified{paper\_conclusion\_hzom\_one\_sided\_contraction\_exponential\_distinguisher\_incompatibility}
设 \((Y,\nu;K_s)\) 为视界零知识算子模型，协议转录满足完美零知识因子化，并且其 Fredholm 行列式
\[
D(s):=\det(I-K_s)
\]
与 \(\Xi_\zeta\) 落在同一解析读出链上。若对每个 \(\varepsilon>0\) 都有
\[
\rho(K_s)<1
\qquad
\bigl(\Re s\ge \tfrac12+\varepsilon\bigr),
\]
则 \(\Xi_\zeta\) 的全部零点都落在临界线 \(\Re s=\tfrac12\) 上。反之，若存在零点
\[
\rho=\beta+\mathrm{i}\gamma,
\qquad
\beta>\tfrac12,
\]
并记其重数为 \(m(\rho)\)，则存在一族有界可见测试函数 \((\Phi_T)_{T\ge 1}\)、真实分布 \((P_T)_{T\ge 1}\) 与常数 \(c>0\)，使任意仅依赖视界因子的模拟分布 \((Q_T)_{T\ge 1}\) 都满足
\[
\bigl|\EE_{P_T}\Phi_T-\EE_{Q_T}\Phi_T\bigr|
\ge
c\,T^{m(\rho)-1}e^{(\beta-\frac12)T}.
\]
因此，单侧谱半径收缩与轴外零点所诱导的指数区分器严格互斥。
\end{theorem}
\begin{proof}
前半部正是主稿中 HZOM 语义下的临界线收缩判据：在任意固定右半平面条带上若都有严格谱收缩，则全部解析谱支撑被压回唯一的临界切片 \(\Re s=\tfrac12\)。

后半部则由同一节的轴外零点放大定理给出。若 \(D(s)\) 在 \(\rho=\beta+\mathrm{i}\gamma\) 处有零点且 \(\beta>\tfrac12\)，则其留数贡献在相应的显式公式中必生成一族有界可见测试 \((\Phi_T)\)，使真实输出与任意仅因子化于视界因子的模拟输出之间出现
\[
T^{m(\rho)-1}e^{(\beta-\frac12)T}
\]
量级的指数偏离。两者合并即得结论。证毕。
\end{proof}

\begin{theorem}[互补零知识通道的临界模平方与谱流化约]\label{thm:conclusion-complementary-zk-channel-critical-modsquare-spectralflow}
设协议在视界层分裂为两条互不通信的互补通道，其行列式因子满足
\[
D(s)=D_+(s)D_-(s),
\qquad
D_-(s)=D_+(1-s).
\]
则在临界切片上有严格实平方律
\[
D\!\left(\tfrac12+\mathrm{i}t\right)
=
\left|D_+\!\left(\tfrac12+\mathrm{i}t\right)\right|^2
\in \RR_{\ge 0}.
\]
若进一步存在有限维酉路径 \(U_+(t)\) 使
\[
D_+\!\left(\tfrac12+\mathrm{i}t\right)=\det\!\bigl(I-U_+(t)\bigr),
\]
则任一临界零点 \(t_0\) 的消失重数恰等于
\[
\dim\ker\!\bigl(I-U_+(t_0)\bigr).
\]
在横截情形下，零点的有向计数等于 \(U_+\) 相对特征值 \(1\) 的谱流。
\end{theorem}
\begin{proof}
在本文互补通道设定中，沿临界切片有 \(1-(\tfrac12+\mathrm{i}t)=\tfrac12-\mathrm{i}t\)，而 determinant 因子满足 Schwarz 反射律，故
\[
D_-\!\left(\tfrac12+\mathrm{i}t\right)
=
D_+\!\left(\tfrac12-\mathrm{i}t\right)
=
\overline{D_+\!\left(\tfrac12+\mathrm{i}t\right)}.
\]
代入 \(D=D_+D_-\) 即得
\[
D\!\left(\tfrac12+\mathrm{i}t\right)
=
\left|D_+\!\left(\tfrac12+\mathrm{i}t\right)\right|^2.
\]

若再有
\(
D_+(\tfrac12+\mathrm{i}t)=\det(I-U_+(t)),
\)
则 \(\det(I-U_+(t_0))\) 的消失重数正等于特征值 \(1\) 的代数重数；在酉情形这又等于
\(
\dim\ker(I-U_+(t_0)).
\)
横截穿越时，有向零点计数与相对特征值 \(1\) 的谱流相同，故总体零点问题被严格化约为单侧酉谱流。证毕。
\end{proof}

\begin{theorem}[Hardy 阻尼临界指数的精确偏移]\label{thm:conclusion-hardy-damping-critical-exponent-exact-shift}
对任意底数 \(b>1\)，定义阻尼 Abel 生成函数
\[
A_{b,\delta}(r):=A_b(b^{-\delta}r),
\]
并令
\[
\delta_H(b):=\inf\{\delta\ge 0:\ A_{b,\delta}\in H^2(\mathbb D)\}.
\]
若记
\[
\sigma_*:=\sup\{\Re\rho:\ \Xi_\zeta(\rho)=0\},
\]
则有精确恒等式
\[
\delta_H(b)=\sigma_*-\frac12.
\]
因而
\[
\mathrm{RH}
\iff
\delta_H(b)=0
\]
对某个 \(b>1\) 成立；等价地，对每个 \(b>1\) 都成立。
\end{theorem}
\begin{proof}
主稿已证明 \(A_{b,\delta}\) 的极点恰位于
\[
r_{\rho,\delta}=b^{-(\rho-(1/2+\delta))},
\]
故“全部极点离开单位圆盘”与
\[
\Re\rho\le \frac12+\delta
\]
完全等价。另一方面，Hardy 门槛判据表明
\[
A_{b,\delta}\in H^2(\mathbb D)
\]
当且仅当存在 \(\varepsilon>0\) 使所有零点都满足
\[
\Re\rho\le \frac12+\delta-\varepsilon.
\]
对满足该条件的 \(\delta\) 取下确界，即得
\[
\delta_H(b)=\sigma_*-\frac12.
\]
于是 \(\delta_H(b)=0\) 与 \(\sigma_*=\tfrac12\) 等价，也即与 RH 等价。证毕。
\end{proof}

\begin{theorem}[Li 第一常数的底数无关能量律]\label{thm:conclusion-li-first-constant-base-independent-energy-law}
在 RH 与零点单纯假设下，定义
\[
E_b(\delta):=\|A_{b,\delta}\|_{H^2(\mathbb D)}^2.
\]
则存在严格重整化极限
\[
\lim_{\delta\downarrow 0}(2\delta\log b)\,E_b(\delta)
=
\sum_{\rho}\frac{1}{\rho(1-\rho)}
=
2+\gamma-\log(4\pi).
\]
因此，临界发散首项的系数与底数 \(b\) 无关，并与 Li 第一系数精确一致。
\end{theorem}
\begin{proof}
主稿已经把 \(H^2\) 能量写成各极点贡献的和，并证明在 RH 与零点单纯条件下，
\[
(2\delta\log b)\,\|A_{b,\delta}\|_{H^2(\mathbb D)}^2
\longrightarrow
\sum_{\rho}\frac{1}{\rho(1-\rho)}.
\]
随后再利用对应的显式恒等式，可把右端化简为
\[
2+\gamma-\log(4\pi).
\]
该表达式不含 \(b\)，故临界发散首项系数为底数无关常数。证毕。
\end{proof}

\begin{theorem}[增长指数、Hardy 门槛与最右零点实部的三重同一]\label{thm:conclusion-growth-hardy-threshold-rightmost-zero-triple-identity}
定义离散显式公式系数的增长指数
\[
\kappa_b:=\limsup_{n\to\infty}\frac{1}{n}\log|E_b(n)|,
\]
并定义能量临界指数
\[
\delta_2(b):=\inf\{\delta\ge 0:\ E_b(\delta)<\infty\}.
\]
则有严格三重同一
\[
\sigma_*-\frac12
=
\frac{\kappa_b}{\log b}
=
\delta_2(b).
\]
因此，最右零点实部偏移、离散系数增长率与 Hardy 能量门槛只是同一个临界量的三种坐标表达。
\end{theorem}
\begin{proof}
主稿分别给出两条恒等式：
\[
\sigma_*=\frac12+\frac{\kappa_b}{\log b},
\qquad
\delta_2(b)=\sigma_*-\frac12.
\]
消去 \(\sigma_*\) 即得
\[
\frac{\kappa_b}{\log b}=\delta_2(b)=\sigma_*-\frac12.
\]
证毕。
\end{proof}

\begin{theorem}[底数幂重整化的反向鞅与 \texorpdfstring{$k$}{k}-进能量分屑]\label{thm:conclusion-basepower-renormalization-reverse-martingale-kadic-splitting}
\leanverified{paper\_conclusion\_basepower\_renormalization\_reverse\_martingale\_kadic\_splitting}
固定整数 \(k\ge 2\) 与 \(\delta>\delta_2(b)\)。定义
\[
F_m(e^{\mathrm{i}\theta}):=A_{b^{k^m},\delta}(e^{\mathrm{i}\theta})
\qquad (m\ge 0).
\]
则边界值族 \((F_m)_{m\ge 0}\) 关于递减过滤
\[
\mathscr F_m:=\sigma(e^{\mathrm{i}k^m\theta})
\]
构成反向鞅，并满足正交能量分解
\[
\|F_0\|_{H^2(\mathbb D)}^2
=
\|F_M\|_{H^2(\mathbb D)}^2
+
\sum_{m=0}^{M-1}\|F_m-F_{m+1}\|_{H^2(\mathbb D)}^2.
\]
更精确地，\(\|F_m-F_{m+1}\|_{H^2(\mathbb D)}^2\) 恰等于 Fourier 系数中 \(k\)-进赋值精确等于 \(m\) 的壳层能量。
\end{theorem}
\begin{proof}
主稿的根单位平均公式给出
\[
A_{b^k,\delta}(r)
=
\frac1k\sum_{j=0}^{k-1}A_{b,\delta}(\omega_k^j r^{1/k}).
\]
取边界值后，该公式正是对 \(\mathscr F_1=\sigma(e^{\mathrm{i}k\theta})\) 的条件期望；迭代即得
\[
F_{m+1}=\EE(F_m\mid \mathscr F_{m+1}),
\]
因而 \((F_m)\) 构成反向鞅。

由于 \(H^2(\mathbb D)\) 的边界值空间是 Hilbert 空间，条件期望差满足勾股恒等式，故对任意 \(M\ge 1\),
\[
\|F_0\|_{H^2(\mathbb D)}^2
=
\|F_M\|_{H^2(\mathbb D)}^2
+
\sum_{m=0}^{M-1}\|F_m-F_{m+1}\|_{H^2(\mathbb D)}^2.
\]
而 \(\mathscr F_m\) 的 Fourier 投影恰对应“频率的 \(k\)-进赋值至少为 \(m\)”的部分，因此差分 \(F_m-F_{m+1}\) 正好抽出 \(k\)-进赋值精确等于 \(m\) 的壳层。其 \(H^2\) 能量便是该壳层的 Parseval 质量。证毕。
\end{proof}

\begin{theorem}[离临界偏移的有限尺度内域显影]\label{thm:conclusion-offcritical-finite-scale-interior-pole-exposure}
设存在零点
\[
\rho=\frac12+\varepsilon+\mathrm{i}\gamma,
\qquad
\varepsilon>0.
\]
取任意 \(0\le \delta<\varepsilon\) 与任意 \(\eta\in(0,1)\)。则只要
\[
m\ge \frac{\log(1/\eta)}{(\varepsilon-\delta)\log b},
\]
阻尼并作 \(m\) 次底数幂重整化后的函数 \(A_{b^m,\delta}\) 必在圆盘 \(|r|\le \eta\) 内出现极点。

因此，任何偏离临界线的零点都能在有限次底数放大后，被显影为任意给定内半径上的显式极点。
\end{theorem}
\begin{proof}
由极点位置公式，
\[
r_{\rho,\delta}(b)=b^{-(\rho-(1/2+\delta))},
\qquad
|r_{\rho,\delta}(b)|=b^{-(\varepsilon-\delta)}.
\]
在底数幂重整化 \(b\mapsto b^m\) 下，极点半径变为
\[
|r_{\rho,\delta}(b^m)|
=
b^{-m(\varepsilon-\delta)}.
\]
若该半径不超过 \(\eta\)，便说明极点已进入圆盘 \(|r|\le \eta\)。解不等式
\[
b^{-m(\varepsilon-\delta)}\le \eta
\]
即得
\[
m\ge \frac{\log(1/\eta)}{(\varepsilon-\delta)\log b}.
\]
证毕。
\end{proof}


\paragraph{组织视角}
上述命题自然分成六个层次：Cesaro--Haar 化及其 \(40\) 状态单通道饱和，把 RH-sharp 临界线梳齿提升为统计序参量；正性 Fredholm 完成中的边界模态簇不灭性给出完成层 obstruction；universal solenoid 的有限见证与共尾有限性构成 finite-witness principle；Artin-visible hull 的共尾固定商不可能性则说明，一切真正保持 determinant 因子与主导类函数见证的终端证书，不可能长期压缩到单一有限商，而必须沿分辨率产生增长且非阿贝尔的可见壳；\(\varphi\)-adic Mellin 指纹的极点格刚性进一步说明有限编译侧只能产生纯虚周期晶格而不允许孤立极点；而 arithmetic solenoid 上的经典 RH 重述，最终把 unitary slice、圆盘零点禁入与 Toeplitz--PSD 层级压缩为同一条解析相位证书链。就此而言，有限见证、完成障碍、可见壳增长、\(\varphi\)-adic 极点晶格与 classical RH 的正性重写在本文语言中不再彼此分离，而被统一到同一条 solenoidal 解析几何接口上。

\endinput
